% ===================================================================
%  সারণি নং 14 — রাস্তা। — দোকানদারেরা।
%
%  Ceci n'est PAS une transcription : c'est une traduction, faite en
%  2026, en bengali d'aujourd'hui. Voir texto/bn/10-sarani-01.tex
%  pour la consigne, qui vaut pour les seize fichiers.
% ===================================================================

%%K t14-apar-1 apar
\VUsaut{4.0mm}
\VUtitre{15.0pt}{60}{সারণি নং 14}
\VUsaut{3.0mm}

%%K t14-tit-1 sub
\VUpk{11.4pt}{40}{রাস্তা। --- দোকানদারেরা।}
\VUsaut{3.0mm}

%%K t14-01-1 p
1. --- আমার বন্ধু জন, যে গ্রামে থাকে, আমার পরিবারের সঙ্গে তিন দিন কাটাতে
এসেছে। \VUgras{এ পর্যন্ত} তার কখনও কোনো বড় নগর দেখার সুযোগ হয়নি, তাই আমরা
নিজেদের সব দিন ঘুরে কাটাই আর যা তার জানা নেই তা আমি তাকে বুঝিয়ে দিই।

%%K t14-02-1 p
2. --- প্রথমে আমরা \VUgras{মানমন্দির} \textsuperscript{(1)} দেখতে গেলাম, যাতে তার
বড় আগ্রহ হল। সেই \VUgras{রাস্তা} \textsuperscript{(3)} দিয়ে ফেরার সময় যেখানে
\VUgras{তুর্কি হাম্মাম} \textsuperscript{(2)}, আমরা একটি \VUgras{অন্ত্যেষ্টির}
\textsuperscript{(4)} মুখোমুখি হলাম। \VUgras{শবযাত্রার} আগে আগে
\VUgras{পুরোহিতেরা} চলছিলেন, সেই \VUgras{শবগাড়ির} আগে যাতে \VUgras{কফিন} রাখা
ছিল। অনেক বন্ধু \VUgras{শোকে} ডোবা পরিবারের পিছনে চলছিলেন।

%%K t14-02-2 p
শোভাযাত্রা পেরিয়ে যেতে আমরা বড় \VUgras{মদের দোকানের} \textsuperscript{(5)}
সামনে রাস্তা পেরোলাম, যেখানে অনেক \VUgras{খদ্দের} ছাদের নিচে \VUgras{বিয়ার}
খাচ্ছিল, কাচের \VUgras{ছাউনির} \textsuperscript{(6)} আড়ালে।

%%K t14-03-1 p
3. --- জন সেই \VUgras{রাজচিহ্ন} দেখে ধন্দে পড়ল যা \VUgras{মদের ব্যবসায়ীর}
\textsuperscript{(7)} দোকান ও \VUgras{সংগীত-বিক্রেতার} \textsuperscript{(10)}
দোকানের মাঝে লাগানো। সে আমাকে জিজ্ঞেস করল তার অর্থ কী; আমি বললাম যে সেটি
ইংরেজ \VUgras{বাণিজ্যদূতাবাসের} \textsuperscript{(8)} পরিচয়, আর তাকে বারান্দায়
দাঁড়ানো \VUgras{বাণিজ্যদূতকে} \textsuperscript{(9)} দেখালাম।

%%K t14-tit-2 sub
\VUpk{10.2pt}{40}{দোকান দেখা।}
\VUsaut{3.0mm}

%%K t14-04-1 p
4. --- স্বভাবতই আমরা \VUgras{বাদ্যযন্ত্র}, \VUgras{হারমোনিয়াম},
\VUgras{অর্গান} ও \VUgras{পিয়ানোর} সাজসজ্জার সামনে থেমে গেলাম;
\VUgras{স্বরলিপির} ও পিয়ানোর \VUgras{খণ্ডের} চিত্রিত মলাট দেখলাম, আর দোকানের
সাইনবোর্ড হিসেবে লাগানো সোনালি কাঠের \VUgras{বীণা} নিহারলাম।

%%K t14-05-1 p
5. --- \VUgras{ঝাড়ুদারদের} \textsuperscript{(11)} ও \VUgras{ঝাড়ুর গাড়ির}
\textsuperscript{(12)} তোলা ধুলোর মেঘ আমাদের বাধ্য করল \VUgras{আসবাব-ব্যবসায়ীর}
\textsuperscript{(13)} সুন্দর \VUgras{আসবাব-সমূহ} ও \VUgras{লোহার ব্যবসায়ীর}
\textsuperscript{(14)} \VUgras{উনুন} ও \VUgras{বাতির} সাজসজ্জার পাশ দিয়ে
তাড়াতাড়ি বেরিয়ে যেতে।

%%K t14-06-1 p
6. --- সেই জায়গাতেই আমাদের ফুটপাত ছাড়তে হল, কারণ তা সেই পুঁটুলিতে আটকে ছিল
যা একটি \VUgras{আসবাব-বওয়া গাড়ি} \textsuperscript{(16)} থেকে নামানো হচ্ছিল,
যাতে সেই \VUgras{দোকানে} \textsuperscript{(15)} নেওয়া হয় যা এইমাত্র ভাড়া
হয়েছে।

%%K t14-06-2 p
এতে আমরা \VUgras{গাড়ি-নির্মাতার} \textsuperscript{(17)} সুন্দর গাড়ি দেখতে
পেলাম না, কিন্তু আমরা থেমে \VUgras{বাঁধনদারকে} \textsuperscript{(19)} এটি দেখতে
থামলাম না যে সে নিজের প্রতিবেশী \VUgras{সাইকেল ও মোটরের ব্যবসায়ীর}
\textsuperscript{(18)} জন্য একটি পেটরা বানাচ্ছে। তারপর, যেহেতু বেশ দেরি হয়ে
গিয়েছিল, আমরা বাড়ি ফিরে এলাম।

%%K t14-tit-3 sub
\VUpk{10.2pt}{40}{নগরে একটি বেড়ানো।}
\VUsaut{3.0mm}

%%K t14-07-1 p
7. --- পরদিন আমার মা প্রস্তাব দিলেন যে জন নিজের ছবি তুলিয়ে নিক, আর আমাকে
বললেন তার সঙ্গে \VUgras{আলোকচিত্রীর} \textsuperscript{(20)} কাছে যেতে। খাওয়ার
পরে আমরা শিল্পীর \VUgras{কক্ষে} গেলাম, যেখানে আমাদের বেশ কিছুক্ষণ অপেক্ষা করতে
হল। সময় কাটাতে আমরা দেয়ালে টাঙানো \VUgras{ছবি} \textsuperscript{(22)}
দেখলাম।

%%K t14-07-2 p
আমাদের পালা এলে কাজে দেরি হল না, আর আমার বন্ধুর \VUgras{ক্যামেরার}
\textsuperscript{(21)} সামনে বসা হয়ে গেলে আমরা নিচে নেমে এলাম।

%%K t14-08-1 p
8. --- প্রথম তলায় আমরা \VUgras{নাপিতের} \textsuperscript{(23)} খোলা দরজা দিয়ে
দেখলাম যে তার সহকারী এক খদ্দেরের চুল কাটছে।

%%K t14-08-2 p
তারপর রাস্তায় এসে আমরা \VUgras{তামাকের দোকানের} \textsuperscript{(24)} পাশ দিয়ে
গেলাম। জনের লাল \VUgras{লণ্ঠন} \textsuperscript{(25)} ও \VUgras{সাইনবোর্ড}
\textsuperscript{(26)} হিসেবে লাগানো \VUgras{বড় কলকেতে} এত মজা লাগল যে সে দুই
বুড়ো ভদ্রলোকের সঙ্গে ধাক্কা খেল, যাঁরা \VUgras{নস্যির চিমটি}
\textsuperscript{(27)} নিতে নিতে কথা বলছিলেন।

%%K t14-tit-4 sub
\VUpk{10.2pt}{40}{মিষ্টির দোকান।}
\VUsaut{3.0mm}

%%K t14-09-1 p
9. --- \VUgras{পোদ্দারের} \textsuperscript{(29)} \VUgras{জানালায়} সাজানো সব
দেশের \VUgras{নোট} ও \VUgras{মুদ্রা} কিছুক্ষণ আমাদের নজর বেঁধে রাখল, কিন্তু
যেহেতু চায়ের সময় হয়েছিল, আমরা আর না থেমে \VUgras{ময়রার}
\textsuperscript{(31)} দোকানে ঢুকে পড়লাম।

%%K t14-10-1 p
10. --- দোকানে যত \VUgras{ভালো জিনিস} ছিল --- নানা রকম \VUgras{কেক,
আদা-রুটি, বিস্কুট} ও \VUgras{মিষ্টি} --- সেসব দেখে আমরা বাছতেই পারলাম না।
শেষে জন \VUgras{সরের কেক} নিল, আর আমি কেবল একটি ছোট \VUgras{চেরির টার্ট},
কারণ মিষ্টি জিনিসে \VUgras{আমার দাঁত ব্যথা করে} আর
\VUgras{দন্তচিকিৎসকের} \textsuperscript{(30)} কাছে বারবার যাওয়ার আমার একটুও
ইচ্ছে নেই।

%%K t14-tit-5 sub
\VUpk{10.2pt}{40}{টাইপ করা।}
\VUsaut{3.0mm}

%%K t14-11-1 p
11. --- নিজের বন্ধুকে একটি \VUgras{টাইপরাইটার} দেখাতে, যার কথা সে শুনেছিল
কিন্তু দেখেনি, আমরা সেই \VUgras{দপ্তরে} \textsuperscript{(32)} গেলাম যা আমার
\VUgras{খুড়তুতো ভাইয়ের} \textsuperscript{(34)}। আমরা তাঁকে পেলাম নিজের
\VUgras{সচিবকে} \textsuperscript{(35)} চিঠি লিখিয়ে দিতে, যিনি
\VUgras{সংকেতলিপিকার}। একটি চিঠি শেষ হতে না হতেই একজন \VUgras{টাইপিস্ট}
\textsuperscript{(33)} সেটি নিজের যন্ত্রে তুলে নিতেন। নিজের আত্মীয়ের কাজ থামাতে
না চেয়ে আমরা তাঁর কাছে কেবল কয়েক মিনিট রইলাম।

%%K t14-12-1 p
12. --- আমার খুড়তুতো ভাইয়ের দপ্তরের নিচে একজন বড় \VUgras{ঘড়িকার ও জহুরি}
\textsuperscript{(37)} থাকেন, যিনি রুপোর কারিগরও। তাঁর চমৎকার দোকানে অনেক
\VUgras{ঘড়ি, অগ্নিকুণ্ডের ঘড়ি, পকেটঘড়ি}, \VUgras{চেন} ও দামি \VUgras{গহনা},
যেমন \VUgras{মুক্তার হার}, সোনার \VUgras{বালা}, \VUgras{কানের দুল} ও
\VUgras{দামি পাথর} এবং \VUgras{হীরে} বসানো \VUgras{আংটি}। একটি জানালা পুরোপুরি
\VUgras{রুপোর বাসনকে} দেওয়া আর তাতে রুপোর \VUgras{চামচ-কাঁটার} একটি বড়
সংগ্রহ।

%%K t14-13-1 p
13. --- রাস্তার মোড় ঘুরতে ঘুরতে আমি দেখলাম যে বাড়ির \VUgras{নম্বরের} পাশে
লাগানো \VUgras{নামফলক} \textsuperscript{(36)} বদলে দেওয়া হয়েছে। আমি এটি জোরে
বলতেই যাচ্ছিলাম, এমন সময় একটি সামরিক ব্যান্ডের প্রসন্ন ধ্বনি শোনা গেল। আমরা
রেজিমেন্টের দিকে দৌড়ে পড়লাম, এটি না ভেবে যে তা \VUgras{টুপিওয়ালার}
\textsuperscript{(39)} ও \VUgras{দস্তানাওয়ালার} \textsuperscript{(40)} দোকানের
সামনে দিয়ে যাবে, আর যে সেটি বেরোতে দেখার জন্য আমরা \VUgras{চশমাওয়ালার}
\textsuperscript{(38)} জানালার সামনে দাঁড়িয়ে থেকেও ততটাই ভালো জায়গায় থাকতাম,
যা \VUgras{নাক-চশমা, চশমা} ও \VUgras{দূরবিনে} ভরা।

%%K t14-tit-6 sub
\VUpk{10.2pt}{40}{সালামি-কুচকাওয়াজ।}
\VUsaut{3.0mm}

%%K t14-14-1 p
14. --- \VUgras{রেজিমেন্টের} \textsuperscript{(41)} আগে আগে \VUgras{ঢোল-প্রধান}
\textsuperscript{(42)} চলছিল, তার পিছনে \VUgras{ঢোলিরা} \textsuperscript{(43)},
\VUgras{বিউগলবাদকেরা} \textsuperscript{(44)} ও সমস্ত \VUgras{ব্যান্ড}।
\VUgras{জেনারেল} \textsuperscript{(45)}, যিনি নিজের সহকারীর সঙ্গে চত্বরে ছিলেন,
\VUgras{জলধারার} \textsuperscript{(49)} পাশে থেমেছিলেন, যা \VUgras{ফোয়ারার}
\textsuperscript{(48)}, \VUgras{কুচকাওয়াজ} দেখতে।

%%K t14-14-2 p
\VUgras{স্টাফের অফিসারেরা} \textsuperscript{(46)}, \VUgras{কর্নেল} ও
\VUgras{উপ-কর্নেল} তাঁকে তলোয়ারে সালাম দিলেন। \VUgras{মেজরেরা} নিজেদের
\VUgras{ব্যাটালিয়নের} \textsuperscript{(47)} আগে ছিলেন আর প্রতিটি
\VUgras{ক্যাপ্টেন} নিজের \VUgras{কোম্পানির} নেতৃত্ব দিচ্ছিলেন, যখন
\VUgras{লেফটেন্যান্ট}, \VUgras{উপ-লেফটেন্যান্ট} ও \VUgras{অ-কমিশনপ্রাপ্ত
অফিসারেরা} নিজেদের লোকদের পাশে নিজেদের সাধারণ জায়গায় ছিলেন।

%%K t14-15-1 p
15. --- \VUgras{চৌরাস্তায়} \textsuperscript{(50)} অনেক পথচারী জমে গিয়েছিল;
দোকানদারেরা --- \VUgras{ছাতা-বিক্রেতা} \textsuperscript{(51)}, \VUgras{রঙরেজ}
\textsuperscript{(53)}, \VUgras{বন্দুককার} \textsuperscript{(54)} ---
\VUgras{সিপাহি} দেখতে বাইরে এল। সব গাড়ি --- \VUgras{ভাড়ার গাড়ি}
\textsuperscript{(52)}, \VUgras{ব্যক্তিগত গাড়ি} \textsuperscript{(57)} ও
\VUgras{জলের গাড়িও} \textsuperscript{(56)} --- ফুটপাতের ধারে দাঁড়িয়ে পড়ল।

%%K t14-tit-7 sub
\VUpk{10.2pt}{40}{রাস্তার দৃশ্য।}
\VUsaut{3.0mm}

%%K t14-15-2 p
হাতে নিজের \VUgras{নমুনার বাক্স} নিয়ে একজন \VUgras{ভ্রাম্যমাণ বিক্রেতা}
\textsuperscript{(55)} তৎপরতায় রাস্তা পেরিয়ে গেল, আর \VUgras{উপরের পাটাতনে}
\textsuperscript{(59)} বসা লোকেরা, যা \VUgras{অমনিবাসের} \textsuperscript{(58)},
ঘুরে তামাশা দেখতে লাগল। এক বেচারা \VUgras{পথের গায়ককে} \textsuperscript{(60)},
যে নিজের \VUgras{বাক্স-বাজনা} \textsuperscript{(61)} নিয়েছিল, নিজের
\VUgras{গান} মাঝপথে ছাড়তে হল, কারণ ঢোলের শব্দ তার গলা চাপা দিয়ে দিল।

%%K t14-16-1 p
16. --- রেজিমেন্ট যখন \VUgras{কাপড়ের দোকানের} \textsuperscript{(64)} সামনে এল,
তখন \VUgras{সেলাইকারিণীরা} \textsuperscript{(63)} ও \VUgras{দরজিরা}
\textsuperscript{(62)} নিজেদের \VUgras{সেলাইয়ের কল} ও নিজেদের কাজ ছেড়ে জানালা
দিয়ে উঁকি দিতে লাগল। \VUgras{দোকানদার} \textsuperscript{(65)}, যিনি এইমাত্র
নিজের \VUgras{সাজানোর জানালায়} নতুন \VUgras{স্যুট} \textsuperscript{(66)}
রেখেছিলেন, তাঁকে \VUgras{কাপড়} \textsuperscript{(67)} ও \VUgras{লিনেনের} সেই
থানগুলি সরাতে হল যা ফুটপাতে \VUgras{নালার মুখের} \textsuperscript{(68)} পাশে রাখা
ছিল, এই ভয়ে যে ভিড় সেগুলি ফেলে না দেয়; এক বুড়ো \VUgras{ফেরিওয়ালা}
\textsuperscript{(69)} পর্যন্ত, যার সঙ্গে এক দারোয়ান কিছু মোজার দর কষছিল, নিজের
বিক্রি সারতে একটি দেউড়িতে ঢুকতে বাধ্য হল।

%%K t14-16-2 p
আমরা রেজিমেন্টের সঙ্গে ব্যারাক পর্যন্ত গেলাম \VUgras{আর}, নিজেদের দিনে খুব
সন্তুষ্ট হয়ে, খাওয়ার সময় বাড়ি ফিরে এলাম।

%%K t14-tit-8 sub
\VUpk{10.2pt}{40}{নানা ব্যবসা ও পেশা।}
\VUsaut{3.0mm}

%%K t14-17-1 p
17. --- পরদিন আমার বাবা প্রস্তাব রাখলেন যে আমাদের স্থানীয় কাগজের ছাপাখানা
দেখিয়ে আনবেন। আমরা উৎসাহে রাজি হলাম।

%%K t14-17-2 p
\VUgras{ছাপাখানার মালিক} \VUgras{পুস্তক-বিক্রেতাও} \textsuperscript{(70)},
\VUgras{প্রকাশক}, \VUgras{লেখার সরঞ্জামের বিক্রেতা} ও \VUgras{বাঁধাইকারও}।
তিনি প্রথম তলায় কাগজের \VUgras{সম্পাদকীয় কক্ষ} \textsuperscript{(71)}
রেখেছেন, আর \VUgras{খোদাই-কক্ষও}, যেখানে \VUgras{শিলামুদ্রকেরা}
\textsuperscript{(72)} ও \VUgras{তাম্র-খোদাইকারেরা} \textsuperscript{(73)} কাজ
করেন, আর শেষে সংযোজন-কক্ষ, যেখানে \VUgras{অক্ষর-সংযোজকেরা}
\textsuperscript{(74)}। আসল \VUgras{ছাপাঘর} \textsuperscript{(75)},
\VUgras{ছাপার কল} ও আলাদা আলাদা যন্ত্র নিচের তলায়।

%%K t14-tit-9 sub
\VUpk{10.2pt}{40}{জন অনেক প্রশ্ন করে।}
\VUsaut{3.0mm}

%%K t14-18-1 p
18. --- ছাপাখানা থেকে বেরোতে বেরোতে জন বড় ধন্দে একটি খুঁটিতে লাগানো ছোট কাচের
বাক্সের সামনে থামল, যা \VUgras{ছোট জিনিসের দোকানের} \textsuperscript{(77)} সামনে
ছিল, আর জিজ্ঞেস করল সেটি কী কাজে লাগে। আমার বাবা বোঝালেন যে সেটি
\VUgras{আগুনের সতর্কবার্তা} \textsuperscript{(76)}। সত্যিই নগরের প্রতিটি জিনিস
আমার বন্ধুর কাছে অবাক করা ছিল, সাধারণ \VUgras{জলের কল} \textsuperscript{(78)} ও
সেই হালকা \VUgras{মালের তিনচাকা} \textsuperscript{(80)} থেকে শুরু করে, যা উর্দি
পরা এক \VUgras{ছোকরা} \textsuperscript{(79)} চালায়, \VUgras{ডাকের গাড়ি}
\textsuperscript{(81)} পর্যন্ত।

%%K t14-19-1 p
19. --- যতক্ষণ আমার বাবা আমাদের ছেড়ে \VUgras{কর-আদায়কারীর}
\textsuperscript{(83)} সঙ্গে কথা বলতে গেলেন, যিনি এক \VUgras{উকিল}
\textsuperscript{(82)} ও এক \VUgras{মুহুরির} \textsuperscript{(84)} সঙ্গে কথা বলতে
থেমেছিলেন, আমরা রাস্তার আনাগোনা দেখে মন বহলাতে থাকলাম। আমাদের সামনে দিয়ে নানা
রকম লোক গেল: এক \VUgras{ময়রার ছোকরা} \textsuperscript{(85)}, যে একটি ঝুড়িতে
একটি চমৎকার \VUgras{পিঠে} নিয়েছিল, এক \VUgras{কাপড়ের ব্যবসায়ীর}
\textsuperscript{(86)} পিছনে এল; এক \VUgras{বাতি জ্বালানেওয়ালা}
\textsuperscript{(87)} এক \VUgras{গ্যাস বসানেওয়ালার} \textsuperscript{(88)} সঙ্গে
কথা বলছিল; এক \VUgras{সন্ন্যাসিনী} \textsuperscript{(89)} এক ছোট্ট অনাথ মেয়ের
হাত ধরে নিজের মঠে ফিরছিলেন।

%%K t14-tit-10 sub
\VUpk{10.2pt}{40}{ফেরার সময়।}
\VUsaut{3.0mm}

%%K t14-20-1 p
20. --- আমার বাবা আমাদের সঙ্গে এসে মিলেছেন কি না, জন দুই \VUgras{চিমনি
পরিষ্কারকের} \textsuperscript{(91)} ঝুলে কালো মুখ ও অদ্ভুত বেশ দেখে হেসে ফেলল।

%%K t14-21-1 p
21. --- আমি নিজের মায়ের জন্য \VUgras{ফুল-বিক্রেতার} \textsuperscript{(92)} কাছ
থেকে একটি গাছ কিনতে চাইছিলাম, কিন্তু আমার বাবা বললেন যে সেটি বইতে খুব ভারী হবে
আর কয়েকটি \VUgras{কমলা} নেওয়াই ভালো। আমরা \VUgras{ঠেলাওয়ালির}
\textsuperscript{(94)} কাছে গেলাম, আর যখন \VUgras{শিক্ষয়িত্রী}
\textsuperscript{(93)}, যিনি নিজের ছোট ছাত্রীর জন্য কয়েকটি বাছছিলেন, নিজের কেনা
শেষ করলেন, তখন আমাদের পালা এল।

%%K t14-22-1 p
22. --- ফেরার পথে আমাদের কয়েকজন চেনা লোকের সঙ্গে দেখা হল: আমার পরিবারের এক
পুরনো সখী, নিজের \VUgras{সঙ্গিনীর} \textsuperscript{(95)} বাহু ধরে, সেতু ও
রাস্তার \VUgras{প্রকৌশলী} \textsuperscript{(96)}, আর আমার এক খুড়তুতো বোন নিজের
\VUgras{ধাত্রীর} \textsuperscript{(98)} সঙ্গে।

%%K t14-22-2 p
তারপর আমরা আমার মায়ের \VUgras{টুপি-নির্মাত্রীর} \textsuperscript{(97)} পাশ
দিয়ে গেলাম আর দুই \VUgras{সন্ন্যাসীর} \textsuperscript{(99)}, যাঁরা নগরে অচেনা
ছিলেন এবং যাঁরা আমার বাবার কাছে স্টেশনের পথ জিজ্ঞেস করতে এলেন।
\VUsaut{6.0mm}
\VUfilet{22mm}
